\chapter{Problem Analysis and Preliminary Study}

\section{System Description}

The system is a networked, multi-tier hospital monitoring solution.
At its base sits the \textbf{ESP32 embedded acquisition node}, which occupies the
sensor layer and embedded layer. Its role is to collect physiological and environmental
measurements, apply local signal processing to derive meaningful values (BPM,
SpO\textsubscript{2}), and transmit structured records to a cloud database.

The system follows a multi-tier architectural pattern designed for high availability and low-latency monitoring. At the edge, the \textbf{IoT Layer} (ESP32) performs continuous data acquisition. This telemetry is pushed to the \textbf{Cloud Persistence Layer} (Firebase), where Realtime Database handles live streams and Firestore manages long-term clinical records (users, people, alerts, and detection events). The \textbf{Service Layer} (Node.js/tRPC) provides a type-safe interface for management operations and AI inference, while the \textbf{Application Layer} (React) delivers a high-fidelity monitoring interface for medical staff. The face-recognition pipeline runs directly in the browser using face-api.js and feeds detection events back to the server for audit logging and alerting.

\begin{figure}[H]
    \centering
	\includegraphics[width=\linewidth]{System Context Diagram.png}
    \caption{System context and operational boundaries.}
    \label{fig:system_context}
\end{figure}

\section{Use Cases}

The following operational scenarios define the expected system behaviour.

\subsection*{Embedded Node Use Cases}

\begin{description}
	\item[UC-1: Initial Configuration.]
	      A technician connects to the ESP32's Access Point and enter Wi-Fi credentials, Firebase credentials, room ID, and patient ID. The device saves all settings to NVS, serves a confirmation page, and reboots.

	\item[UC-2: Normal Operation.]
	      After boot, the device loads credentials from NVS, connects to the hospital Wi-Fi,
	      authenticates with Firebase, and enters continuous monitoring mode.

	\item[UC-3: Wi-Fi Recovery.]
	      If Wi-Fi is lost during operation (or fails on cold start), the embedded device
	      detects the disconnection and retries \texttt{wifiConnect()} every 10 seconds,
	      indefinitely, without manual intervention.

	\item[UC-4: Remote Reconfiguration.]
	      A technician connects to the AP at any time to update credentials.
	      After saving, the device reboots and applies the new configuration.

	\item[UC-5: Sensor Fault Detection.]
	      If a sensor returns an invalid reading (NaN from DHT22, or an IR raw value
	      below the 50\,000-count finger-presence threshold from the MAX30102), the
	      firmware discards the sample, skips the upload, and continues monitoring
	      without writing corrupt data to the database.
\end{description}

\subsection*{Backend Use Cases}

\begin{description}
	\item[UC-6: Staff Authentication and RBAC.]
	      Medical staff (Doctors and Nurses) authenticate via Firebase. The system enforces Role-Based Access Control (RBAC), ensuring that clinical records are only accessible to authorized personnel assigned to the specific patient.
	      
	\item[UC-7: Clinical Data \mbox{Synchronization}.]
	      Administrative users trigger \mbox{synchronization} between the local cache and the Firebase Firestore backend to ensure data consistency across multiple monitoring stations.
	      
	\item[UC-8: Automated Alert Dispatch.]
		      The dashboard applies configurable thresholds to live telemetry; when a violation is detected, it triggers email alerts to the assigned Doctor and Nurse through the backend's SMTP relay. The server also creates alert logs for abnormal presence events (unknown person, patient absence).
\end{description}

\subsection*{Web Dashboard Use Cases}

\begin{description}
	\item[UC-9: Real-time Telemetry Visualization.]
	      Staff view live vitals (BPM, SpO2) and room conditions via high-frequency charts. The dashboard subscribes directly to the Realtime Database to ensure sub-second latency.
	      
	\item[UC-10: Personnel Tracking via Face Recognition.]
		      Using the terminal's camera, the dashboard identifies staff members entering the room and logs their presence. It cross-references detected faces against the authorized staff directory and flags unauthorized entries, applying a multi-frame stability requirement before issuing alerts.
\end{description}

\subsection*{AI Subsystem Use Cases}

\begin{description}
	\item[UC-11: Clinical Anomaly Detection.]
	      The AI layer analyzes time-series windows of patient vitals to identify statistical outliers (Z-score analysis) that may indicate sensor failure or acute physiological changes.
	      
	\item[UC-12: Predictive Vital Forecasting.]
		      Using linear trend analysis, the AI subsystem projects the patient's vitals for the subsequent 5 samples (e.g., \mbox{$\approx$5\,s} at a 1\,Hz input rate), providing clinicians with a short-horizon early-warning signal.
\end{description}

\section{Inputs and Outputs}

\subsection*{Embedded Node — Inputs}

\begin{itemize}
	\item \textbf{DHT22 (GPIO~4, one-wire):} digital temperature (°C) and relative
	      humidity (\%).
	\item \textbf{MAX30102 (I\textsuperscript{2}C fast mode):} raw infrared
	      and red photodiode FIFO values, used to compute BPM and SpO\textsubscript{2}.
\end{itemize}

\subsection*{Embedded Node — Outputs}

\begin{itemize}
	\item \textbf{Firebase RTDB — \texttt{/rooms/\{roomId\}}:}
	      JSON records containing \texttt{temperature}, \texttt{humidity}, and
	      server-side \texttt{timestamp}, pushed on significant change.
	\item \textbf{Firebase RTDB — \texttt{/patients/\{patientId\}}:}
	      JSON records containing \texttt{heartRate}, \texttt{spO2}, and
	      server-side \texttt{timestamp}, pushed on significant change.
\end{itemize}

	\noindent\textbf{Identifier alignment.} The dashboard subscribes to
	\texttt{/patients/\{firebaseId\}} for the selected patient. The ESP32 configuration
	therefore uses the patient's Firebase identifier as \texttt{patientId} to ensure
	consistent stream access across subsystems.

\subsection*{Backend, Web Dashboard, and AI — Inputs/Outputs}

\begin{itemize}
	\item \textbf{Backend Inputs:} tRPC procedure calls (JSON), Firebase ID Tokens, and Firestore/RTDB document snapshots.
	\item \textbf{Backend Outputs:} Structured clinical records (Users, People), alert logs, detection events, and SMTP-formatted email alerts when invoked by the dashboard.
	\item \textbf{Web Inputs:} Firebase RTDB telemetry stream, face-api.js descriptor vectors, and tRPC query results.
	\item \textbf{Web Outputs:} Real-time SVG/Canvas visualizations, detection event logs, and automated toast notifications.
	\item \textbf{AI Inputs:} Sliding-window vectors of Heart Rate and SpO2 values (target length 20; minimum 5 required for anomaly detection).
	\item \textbf{AI Outputs:} Rule-based clinical status (stable/warning/critical), anomaly flags, and 5-step trend forecasts.
\end{itemize}

\section{Constraints}

\subsection*{Memory Constraints}

The ESP32 has 520\,KiB of internal SRAM.
The Firebase FreeRTOS task is allocated 8\,192 bytes of stack, which must accommodate
the TLS handshake context, async client buffers, and local variables.
The NVS partition stores up to ten string key-value pairs for device configuration.

\subsection*{Latency Constraints}

The Firebase upload interval is \texttt{FIREBASE\_INTERVAL\_MS = 1\,000\,ms},
providing a near-real-time update rate.
The DHT22 enforces a hardware minimum inter-sample interval of 2\,000\,ms.


\subsection*{Energy Constraints}

The device is assumed to be mains-powered in a fixed hospital room installation.
No active power management (light sleep, deep sleep, or radio duty cycling) is
implemented in the current version.

\subsection*{ESP32 Platform Limitations}

\textbf{Sensor timing sensitivity.}
The MAX30102 heart rate and SpO\textsubscript{2} sensor produces a continuous stream
of readings that must be drained from its internal buffer regularly.
If the program is busy doing something else — such as waiting for a Firebase upload
to complete, which can take several hundred milliseconds or more — the sensor buffer
fills up and new samples are dropped.
This results in the sensor reporting zero or no data, as if no finger were present,
even when one is.

\textbf{Why two cores are necessary.}
The ESP32 contains two independent processing cores (Protocol CPU and Application CPU) that can run tasks simultaneously. This hardware feature is the direct solution to the timing conflict described above. By assigning high-priority sensor acquisition exclusively to Core~1 and network-bound Firebase communication exclusively to Core~0, the system achieves \textbf{temporal isolation}. The sensor continues reading at its required frequency (draining the MAX30102 FIFO every 25\,ms) regardless of how long the TLS handshake or Firebase upload takes. This ensures that no cardiac samples are dropped during network congestion, maintaining the integrity of the clinical data stream. Furthermore, Core~0's Wi-Fi management remains responsive, allowing for background SSID scanning and reconnection without introducing jitter into the physiological measurements.
\subsection*{AI and Cloud Constraints}

The AI layer operates under \textbf{latency constraints} where inference must complete within the dashboard refresh cycle (15\,s refetch for AI insights) and should remain well below the 1\,Hz telemetry cadence. Computation is performed on the service layer to preserve the ESP32's limited SRAM for data acquisition. The anomaly model requires a minimum of 5 data points to compute a stable standard deviation.

\subsection*{Client-Side Compute Constraints}

The face recognition pipeline executes in the browser and loads multiple
pre-trained face-api.js models from a CDN. Model load time depends on bandwidth
and browser caching; inference performance depends on WebGL availability and the
client CPU/GPU. The detection loop is rate-limited to one frame every 600\,ms to
avoid saturating the UI thread on mid-range devices.

\subsection*{Data Model Constraints}

The Firestore schema uses collections \texttt{users}, \texttt{people},
\texttt{detectionEvents}, \texttt{alertLogs}, and \texttt{roomActivityLogs}.
The Realtime Database stores sensor streams under \texttt{/patients/\{patientId\}}
and \texttt{/rooms/\{roomId\}}, and mirrors detection/alert logs under
\texttt{/detectionEvents/\{userId\}} and \texttt{/alertLogs/\{userId\}} for
real-time client subscription. A base64 storage fallback is used for face
photos when Firebase Storage is not available.
